\documentclass[11pt,a4paper]{ctexart}

% ------------------------------------------------------------------
% 字体与排版（跨平台优先：Windows 使用中易，macOS/Unicode 回退）
% ------------------------------------------------------------------
\usepackage{iftex}
\ifXeTeX
  \usepackage{fontspec}
  \IfFontExistsTF{SimSun}{
    \setCJKmainfont{SimSun}[BoldFont=SimHei,ItalicFont=KaiTi]
    \setCJKsansfont{SimHei}
    \setCJKmonofont{FangSong}
    \setmainfont{Times New Roman}
    \setsansfont{Arial}
    \setmonofont{Consolas}
  }{
    \setCJKmainfont{Noto Serif CJK SC}
    \setCJKsansfont{Noto Sans CJK SC}
  }
\fi

% ------------------------------------------------------------------
% 版式与宏包
% ------------------------------------------------------------------
\usepackage[a4paper,margin=2.5cm,top=2.8cm,bottom=2.8cm]{geometry}
\usepackage{amsmath,amssymb}
\usepackage{booktabs}
\usepackage{array}
\usepackage{longtable}
\usepackage{multirow}
\usepackage{enumitem}
\usepackage[colorlinks=true,linkcolor=black,citecolor=teal,urlcolor=blue]{hyperref}
\usepackage{microtype}
\usepackage{caption}
\captionsetup{font=small,labelfont=bf,justification=raggedright,singlelinecheck=false}

% 绘图（pgfplots + TikZ）
\usepackage{pgfplots}
\pgfplotsset{compat=1.17}
\usepackage{tikz}
\usetikzlibrary{patterns,positioning,arrows.meta}
\tikzset{every node/.append style={font=\small}}

% ------------------------------------------------------------------
% 标题配色：统一使用沉稳的深蓝，避免零散彩色
% ------------------------------------------------------------------
\usepackage{xcolor}
\definecolor{thcolor}{RGB}{20,62,99}   % 主标题深蓝
\definecolor{htcolor}{RGB}{31,78,121}  % 章节标题蓝
\definecolor{ruled}{RGB}{200,60,40}    % 参考文献编号色（极少用）

\usepackage{titlesec}
\titleformat{\section}
  {\Large\bfseries\color{thcolor}}
  {\thesection}{0.6em}{}
\titlespacing{\section}{0pt}{1.2em}{0.5em}
\titleformat{\subsection}
  {\large\bfseries\color{htcolor}}
  {\thesubsection}{0.6em}{}
\titlespacing{\subsection}{0pt}{0.9em}{0.3em}
\titleformat{\subsubsection}
  {\normalsize\bfseries}
  {\thesubsubsection}{0.5em}{}
\titlespacing{\subsubsection}{0pt}{0.6em}{0.2em}

% ------------------------------------------------------------------
% 表头样式
% ------------------------------------------------------------------
\renewcommand{\arraystretch}{1.25}

\newenvironment{vtable}[1]
  {\begin{longtable}{#1}\toprule}
  {\bottomrule\end{longtable}}

% 每段首行缩进 2 字符，中文排版标准
\usepackage{indentfirst}
\setlength{\parindent}{2em}

\begin{document}

% ==================================================================
% 标题（默认不显示日期；按正文惯例，文档不以生成日期落款）
% ==================================================================
\title{\color{thcolor}\Large\bfseries 波动率即资产：VIX 研究二十五年综述}
\author{量化投研学习笔记}
\date{}
\maketitle

\begin{abstract}
过去二十多年间，芝加哥期权交易所（Cboe）的 VIX 指数从一种概括市场恐慌的抽象指标，逐步发展为拥有期货、期权、ETF/ETN、互换与策略指数的完整资产类别。围绕它的研究也随之从零散的观察，聚合成横跨资产定价、金融计量、市场微观结构、行为金融、宏观金融与机器学习的系统领域。本文按若干研究课题组织评述：VIX 的信息含量与预测能力、波动率反馈与杠杆效应、方差风险溢价、波动率期限结构、衍生品定价与微观结构、交易策略的实证证据、情绪与行为层面、跨国与跨资产联动、极端事件与尾部风险，以及计量边界与机器学习应用。全文以约 2005–2026 年为主要叙事区间，个别课题因需追溯构造源头而适当回溯。文中的里程碑文献、事件日期与指数极值均已先经联网核验；凡属行业常识性数值而非权威一手来源者，一律以区间表述并标注“需人工复核”，不作虚构。
\end{abstract}

\tableofcontents
\clearpage

% ==================================================================
% 引言：综述的对象与边界
% ==================================================================
\section*{引言：这篇综述回答什么问题}
\addcontentsline{toc}{section}{引言}

期权价格同时包含两样东西：对波动本身的预期，以及人们为对冲这种波动愿意支付的溢价。VIX 的价值在于把这两者压缩进一个可由期权组合静态复制的数值——这使它既能被观测，也能被交易。本文据此把 VIX 视作研究“不确定性、风险厌恶与情绪”三者关系的实验室，也视作一种携带自身期限结构、偏度与极端形态的风险资产。

在进入具体课题之前，有必要先明确 VIX 的定义与计算方法。VIX 全称芝加哥期权交易所（Cboe）波动率指数，是标准普尔 500 指数（S\&P 500）未来三十天隐含波动率的市场定价，以年化的百分数表示：一个 18 的读数，可解读为期权市场预期标普 500 在未来一年内约有一个标准差、即约 18\% 的波动幅度。它并非由某一个标的价格直接可交易，而是从一整组 S\&P 500 指数期权的价格中反解出的“隐含波动率”。

计算上，VIX 采用模型自由（model-free）的方差加权方法，这在 2003 年被 Cboe 采用并延续至今。其核心思想是：一段区间的已实现方差，可以仅用一篮子虚值（out-of-the-money）看涨与看跌期权的加权价格来复制，因此不必指定任何特定的随机过程，也不必依赖 Black-Scholes 式的单一平值隐含波动率。给定某个执行价 $K$ 上期权的价格，VIX 公式大致为
\[
\sigma^2 \;=\; \frac{2}{T}\sum_i \frac{\Delta K_i}{K_i^2}\;e^{rT}\,O(K_i) \;-\; \frac{1}{T}\Big(\frac{F}{K_0}-1\Big)^2,
\]
其中 $T$ 为剩余期限（年），$K_i$ 为离散的执行价，$\Delta K_i$ 为相邻执行价间隔，$O(K_i)$ 为执行价 $K_i$ 处虚值期权的买卖中间价，$r$ 为无风险利率，$F$ 为远期指数价格，$K_0$ 为低于远期价的最后一个执行价。公式对多个月份的期权进行加权插值，使到期期限恰为三十天；对求得的方差开平方并年化，即得到 VIX。这一计算可回溯至 1990 年。需要强调，模型自由指它不依赖特定参数，但并非构造无干预——它仍受执行价网格疏密、报价质量与数据来源的影响。

更进一步的公式化理解，是把 $VIX^2$ 看作风险中性测度下对三十天方差的风险中性期望（记作 $E^Q[RV]$），而将市场真正发生的已实现方差视为现实测度（$P$）下的期望。二者之差即方差风险溢价：$VRP = E^Q[RV] - E^P[RV]$。当 $VIX^2$ 系统性高于已实现方差时（长期正是如此），期权市场在对波动率收取一个正的溢价，这正是后续各章反复围绕的核心量。

全书贯穿一条统一的分析主轴，值得在这篇综述伊始就点明：\textbf{把 VIX 的一切现象都还原为“风险中性测度与现实测度的差距”（即 $Q$ 与 $P$ 之差）在不同维度上的投影。} 这一差距在波动率层面是方差风险溢价（第四章），在其期限结构层面是逐期限溢价的斜率与“远期波动近乎免费”（第五章），在其偏度层面是负偏的价格（第七、八章），在其极端事件层面是厚尾与跳空（第十章）。读者可把 $Q$ 与 $P$ 的差距想象成一个“价格尺”——它同时标定市场愿意为不确定性支付的拥挤程度，也标定下跌风险的对冲成本。四个设定问题——哪些维度上 $Q$ 与 $P$ 分离、分离的幅度、为何持续存在、以及可被何种结构与机械行为放大——构成理解 VIX 研究的四个根问题，本文后续各章分别围绕它们展开。要在这一框架下读完全文，也需要认识到一条反复出现的经验事实：$Q$ 对下跌与跳空的定价几乎总高于 $P$ 所对应的可实现对冲概率，这正是“保险价格高于精算成本”的镜像，也是“卖波动率长期为正收益”的根本来源。

写作目的有两层。一是把二十多年来的研究主题拆解成若干可独立讨论的课题，理清每个课题的文献脉络与结论边界；二是在“高收益神话”广泛传播的背景下，让读者养成用不对称性来看待这类资产的意识——大多数交易日里盈利微薄的策略，与少数交易日里损失惨重的策略，往往是同一回事。文末附有关键事件年表、术语速查与研究者自查清单，便于检索与复现。

需要事先说明三类陈述的区分。其一，可追溯事实，文中给出文献出处或事件时间；其二，经济学机制的推断，文中给出推理链条；其三，实践判断，文中会标明属判断性结论并给出适用边界。凡标“需人工复核”的数值，均建议以交易所或论文一手数据核验后再引用。

% ==================================================================
% 第一章
% ==================================================================
\section{从指标到资产：VIX 的构造与历史更重要}
% 导语：构造哲学

理解关于 VIX 的研究，先要回答一个问题：一个由期权价格计算出来的数，如何最终成为可以交易、可以对冲、又会让人亏掉身家的资产。答案不在数理技巧里，而在更深的结构——正是因为 VIX 是可静态复制的波动率敞口，它才跨过了“可观测指标”与“可交易商品”之间的门槛。

VIX 的构造刻意回避了“哪一个期权定价模型是正确的”这一争论。1993 年，Cboe 邀请 Robert Whaley 基于标普 100（OEX）平值期权隐含波动率，编制最初的 VIX。这种早期做法依赖 Black-Scholes 世界假设，只用平值一档，对模型误设定敏感。2002 年，高盛交易员 Shah 与 Rattray 重写公式，把思想来源锚定到方差互换定价：一段区间的已实现方差，可以仅用一篮子虚值看涨与看跌期权的加权价格静态复制（Britten-Jones 与 Neuberger 2000 给出了严格条件）。既然如此，“期权市场对 30 天方差的要价”就是一个不依赖具体参数的量。

2003 年，Cboe 与高盛协作推出新版 VIX：改用标普 500（SPX）、覆盖一段执行价区间上的大量虚值期权，采用模型自由（model-free）的方差加权，并回溯至 1990 年的历史数据；最初的标普 100 版本更名 VXO。这一公式“附带”了一条可静态复制的对冲路径，也就建成了让抽象指标变成可交易商品的桥——这构成 VIX 与其余宏观指标的分野：GDP 只能被观察，方差可以被买卖。

随后的产品化节点清晰：2004 年 3 月 24 日 Cboe 期货交易所（CFE）推出 VIX 期货；2006 年 2 月推出 VIX 期权；此后“波动率的波动率”VVIX、9 日 VIX9D，以及原油（OVX）、黄金（GVZ）等跨资产波动率指数陆续登场。2014 年 10 月 6 日 Cboe 引入周度（weekly）SPX 期权以贴近 30 日目标期限，Andersen、Bondarenko 与 Gonzalez-Perez 的测算表明这使到期期限插值误差下降约一个量级；代价是 2014 年前后的 VIX 序列存在统计性质断裂，任何跨断点的实证都应分段或加哑变量处理。

几点对研究者的提醒：新构造可复制，因而衍生出了互换、期货、期权与 ETP 的完整产品线；方差分布的负偏、尖峰与右偏形态决定了它不适于简单正态假设；2003 与 2014 两次改动是两个需显式处理的断点。“模型自由”不等于构造无干预——它仍受执行价网格疏密、数据来源（Cboe 以 C1 交易所单一数据源）、以及深度虚值期权报价质量的影响，极端行情中误差会被放大。简言之，VIX 不是铅那样纯净的量。

\begin{figure}[htbp]
\centering
\begin{tikzpicture}
\begin{axis}[
    width=0.99\textwidth, height=8.0cm,
    xlabel={年份}, ylabel={VIX（收盘）},
    xmin=1989.5, xmax=2030.5, ymin=0, ymax=100,
    xtick={1990,1993,1995,1998,2000,2003,2005,2008,2010,2013,2015,2018,2020,2024,2026},
    xticklabels={1990,1993,1995,1998,2000,2003,2005,2008,2010,2013,2015,2018,2020,2024,2026},
    ytick={0,20,40,60,80,100},
    grid=major, grid style={gray!25},
    axis lines=left,
    enlarge x limits=0, enlarge y limits=0.06,
    legend style={at={(0.02,0.98)},anchor=north west,font=\footnotesize,draw=none},
]
% 关键点位数据（收盘，来源：Cboe 官方与上文核验）；深蓝主曲线
\addplot[thick,thcolor,mark=*,mark size=1.8pt] coordinates {
    (1990,17.6) (1993,12.5) (1995,12.2) (1997,27.2) (1998,43.7) (2000,24.9)
    (2002,27.6) (2003,21.1) (2004,14.7) (2005,12.1) (2006,14.4)
    (2007,22.5) (2008,40.0) (2008.88,80.86) (2010,17.5) (2011,23.4)
    (2014,15.9) (2015,18.2) (2017.84,9.14) (2018.09,37.3) (2018,23.4)
    (2020.21,82.69) (2020,22.8) (2021,17.4) (2022,27.2) (2023,15.5)
    (2024.59,38.6) (2024,15.8) (2025,17.3) (2026,16.5)
};
% 事件点用红色、加大
\addplot[only marks,mark=*,mark options={color=ruled,scale=2.6}] coordinates {(2008.88,80.86) (2020.21,82.69) (2017.84,9.14) (2018.09,37.3) (2024.59,38.6)};
% 极值参考虚线
\draw[dashed,gray] (2008.88,0)--(2008.88,80.86);
\draw[dashed,gray] (2020.21,0)--(2020.21,82.69);
\draw[dashed,gray] (2017.84,0)--(2017.84,9.14);
% 事件标注：三个极值用引线指向不同空白区，避免文字互相重叠
\draw[->,ruled] (axis cs:2008.88,80.86) -- (axis cs:2011.2,90.5) node[anchor=west,fill=white,text=ruled] {\footnotesize 2008 高点 80.86};
\draw[->,ruled] (axis cs:2020.21,82.69) -- (axis cs:2023.2,88.0) node[anchor=west,fill=white,text=ruled] {\footnotesize 2020 收 82.69};
\draw[->,ruled] (axis cs:2017.84,9.14) -- (axis cs:2015.2,17.5) node[anchor=east,fill=white,text=ruled] {\footnotesize 2017 低点 9.14};
\draw[->,ruled] (axis cs:2018.09,37.3) -- (axis cs:2019.9,30) node[anchor=west,fill=white,text=ruled] {\footnotesize 2018 单日 +115\%};
\draw[->,ruled] (axis cs:2024.59,38.6) -- (axis cs:2023.0,50.0) node[anchor=west,fill=white,text=ruled] {\footnotesize 2024 套利平仓};
\end{axis}
\end{tikzpicture}
\caption{VIX 关键点位与事件（1990--2026）。曲线连接代表性年度收盘水平，红线标明已核验的真实极值：2008-11-20 收 80.86、2020-03-16 收 82.69、2017-11-03 收 9.14；2018 与 2024 为年代际事件示意。\\除标注点位外，曲线形态属示意，不作为精确逐日序列引用。}
\label{fig:vixhist}
\end{figure}

% ==================================================================
% 第二章
% ==================================================================
\section{信息含量与预测能力：隐含波动率多知道什么}

研究信息含量，要把 VIX 与已实现波动率的关系摆正。已实现波动率（RV）由真实价格路径计算，属现实测度（P）下的历史事实；期权隐含波动率（VIX）是风险中性测度（Q）下对未来方差的市场定价。二者的差，即方差的 Q–P 差，理论上恰是方差风险溢价。也就是说，即便期权市场理性且无溢价，VIX 也只是对 RV 的最优无偏预测；任何系统性偏差都指向预期之外的成分，如风险厌恶、稀缺性、情绪或预测失误。这条分解是信息含量与方差风险溢价两个主题共用的钥匙。

文献层面的起点是 Whaley（2000）在《Journal of Portfolio Management》发表的《Stock Market Return and the Fear Gauge》，他用“恐慌指标”命名 VIX，并系统展示 VIX 与标普 500 的强负相关，确立了“VIX 上行即风险偏好评判收缩”的叙述基调。思想源头更早在 Brenner 与 Galai（1989）提出的“西格玛指数”。建模一侧，Bollerslev（1986）的 GARCH 建立了条件方差的基准测度，长期作为 VIX 的现实测度对照；Carr 与 Wu（2009）系统检验了隐含—已实现波动率溢价的可持续性。结构性证据由 Bekaert 与 Hoerova（2014）给出，他们主张把 VIX 平方分解为标的条件方差与方差风险溢价两部分，并发现溢价部分预测股票未来收益，而条件方差部分预测真实经济活动——这否定了 VIX 单个数能同时完成两类预测的直觉。他们同时指出 BTZ（2009）把条件方差近似为鞅的假设与数据不符，会带来有偏的溢价估计，由此引出该领域持续的测量之争。

对应用的意义在于，若“溢价”成分确有预测力，择时就不应只看 VIX 绝对水平，而应看 VIX 平方相对条件方差的高估程度；高阶的 VIX 也未必意谓即将崩跌，而可能只是对冲需求高企的定价，这两层含义对期权买卖双方并不相同。

需要说明这条研究线的局限。条件方差由哪种模型估计，本身就会改变结论（所谓“模型炸弹”）；同时可解释的 R 平方很小，统计显著与经济上可利用并不等同。近年文献（如 Lochstoer 与 Muir）借助调查数据与远期期权，指出市场对波动率的预期是缓慢移动的，这提示 VIX 捕捉的或许不是原子式的理性预期，而是具有粘性的预期修正。

% ==================================================================
% 第三章
% ==================================================================
\section{波动率反馈、杠杆效应与风险—收益权衡}

学术界对股票市场存在一个长期反直觉的观察：波动率与当期收益强负相关，坏消息往往放大波动；但“更高的波动率是否应带来更高的预期收益”，却一直缺乏有力支持。VIX 为观察这个风险—收益悖论提供了直接载体。

两种互为镜像的机制常被并置讨论。杠杆效应（leverage effect）是后向的：股价下跌使权益杠杆上升，进而抬升波动率——波动的升高是此前价格下跌的结果。波动率反馈（volatility feedback）是前向的：波动率上升会推高风险溢价，进而抬高贴现率，压低头寸的当期价格——波动的升高成了影响价格的原因。区分的核心在于“前瞻预期”与“历史落地”的方向。

从本书的统一框架看，杠杆与反馈分别对应 $Q$ 与 $P$ 两层的不同机制。杠杆效应发生在物理测度 $P$：它描述“已实现价格下跌被当成波动上升的原因”，本质是公司资本结构变化的后向传导，几乎不依赖风险定价。而波动率反馈则发生在风险中性测度 $Q$ 到价格的反向传导：它要求“波动率上升 → 需要更高的风险溢价 → 贴现率上升”，即 $Q$ 对风险定价的调整反作用于价格。正因为两条通道分属两个测度，才可以在观测上把它们区分开——若只作 $P$ 测度下的统计（收益与波动的同期相关），很难截住反馈通道；要识别它，通常需借助条件波动率对未来收益预测的方向性（杠杆为后向、反馈为前向）。这也解释了为何“波动率上升是否抬高预期收益”在实证里如此不稳：若波动率的上升主要经由反馈（前向贴现率）而非杠杆（后向结构）发生，则当期收益被压、而未来收益尚不确定，风险—收益权衡在不同样本期便会出现方向相反的证据。

经典证据来自 French、Schwert 与 Stambaugh（1987），他们发现已实现收益与已实现波动率的当期创新强负相关，这更像贴现率效应的产物；而 Glosten、Jagannathan 与 Runkle（1993）等则发现条件波动率对预期收益的解释力很弱甚至为负。Lochstoer 与 Muir（2019）用调查与远期期权证明市场对波动率的预期缓慢移动，先对波动消息反应不足、随后补涨，从而在较短期限呈弱风险—收益权衡，在较长期限呈较强权衡。Moreira 与 Muir（2017）则把这种弱关联转化为波动率择时的来源：在已实现波动较低时加仓、较高时减仓，样本内可显著提升风险调整后的绩效。

对趋势或动量交易者而言，波动率反馈意味着高波动往往压制当期价格，这解释了在低波动环境下建仓的头寸，如何在波动激增时同时受到价格与波动的双重打击。更一般的含义是把“杠杆分散”与“时间分散”重新权衡——按波动率反比配置的仓位，本质上是在用时间维度的分散替代一层标的维度的分散。

局限在于，反馈与杠杆在观察上高度共线，很难彻底分离；样本内择时收益对费率、执行与波动率估计极为敏感，回测的优越性在离开无摩擦、可卖空等理想化假设后会明显衰减（具体口径需人工复核）。

% ==================================================================
% 第四章
% ==================================================================
\section{方差风险溢价：被反复检验的预测之谜}

方差风险溢价（Variance Risk Premium, VRP）问的是最具体的问题：为什么卖出波动率（卖方差互换、卖 VIX 期货或卖出期权）在历史上多数时候赚小钱，却在少数日子里亏掉大钱；这个确定性的正收益究竟是承担了真实风险，还是市场在系统性定价失误。

VRP 的定义本身清晰：风险中性测度下期望的方差，减去现实测度下期望的方差。以数值表述，即 VIX 平方对应的 Q 测度方差，减去同一时点已实现（预期）方差。当 Q 测度方差大于 P 测度方差时，期权或方差互换被系统高估价，卖出方收取正溢价；经验上 VRP 长期为正，这正是卖波动率策略的共同底色。对 VRP 长期为正的经济学解释，主流看法落在风险厌恶与保险需求上：投资者愿意为“波动率上升、尤其是下跌端的波动”支付额外保费，正如住户愿为概率很低但损失巨大的火灾买保险。因此 VRP 既可理解为恐惧的定价，也可视作重尾负偏的镜像。

文献脉络是这一领域最重要的主干。Bollerslev、Tauchen 与 Zhou（2009，RFS 22(11) 4463–4492）是标志性论文，他们以 VIX 的平方（Q 测度方差）与样本期已实现方差之差作为 VRP 的代理，并证明这一溢价可显著预测下个月的标普 500 超额收益——高溢价之后收益倾向走高。需要注意，BTZ 采用期望方差近似为等价鞅的简化，Bekaert 与 Hoerova（2014）随后指出这在某些时期会带来有偏甚至为负的溢价估计。Bekaert 与 Hoerova（2014，JFE）系统梳理多种 VRP 测度，主张把 VIX 平方拆成标的条件方差与方差风险溢价，并发现 VRP 预测股票未来收益、条件方差则预测真实经济活动；Bekaert、Hoerova 与 Lo Duca（2012）进一步以 VAR 揭示货币政策松紧会显著影响 VRP 进而影响风险厌恶，从而把卖波动率的收益部分解释为央行压低均衡风险厌恶的结果。Drechsler 与 Yaron（2011）主张 VRP 与“长风险”（long-run risks）一致，可在一般均衡的随机不确定性下复现。Cheng（2018，RFS，The VIX Premium）转向 VIX 期货，发现“风险上行时 VIX 溢价反而回落”这一看似矛盾的现象，并论证它来自真实的对冲需求下降而非测度错误；其 VIX 溢价的“前向”（ex ante）成分能近乎一比一地预测“后向”（ex post）回报，提示市场定价中存在可利用的系统性预期偏差。Londono 与 Xu（2019）把 VRP 拆成下行（DVP）与上行（UVP）两成分，发现 DVP 为正且逆周期，能对四至六月收益做驼峰形预测，UVP 则只对极短期限有效——这强调了负偏才是溢价的来源。

\textbf{测度转换的机制：VRP 为何能在定价中持续存在。} 要理解 VRP 为什么长期为正，需要推敲它对应的资产组合。在一只标准单名或指数方差互换中，合约在到期时支付“已实现方差减去事先固定的波动率执行价”。设标的价 $S_t$ 在现实测度 $P$ 下满足
\[
\frac{dS_t}{S_t}=\mu_t\,dt+\sigma_t\,dW^P_t ,
\]
则在无套利定价下，波动率执行价恰等于风险中性测度 $Q$ 对方差的期望，而该值一般不等于现实测度 $P$ 下的条件方差。用伊藤引理并结合随机测度变换，可把 VRP 写成两种测度下期权复权积分的差；这里给出一个更可操作的等价形式。设利率为零、风险中性测度 $Q$ 与物理测度 $P$ 之间的密度 $dQ/dP=M_T$ 存在，则
\[
E^Q\!\Big[\int_0^T\sigma_t^2\,dt\Big]\;-\;E^P\!\Big[\int_0^T\sigma_t^2\,dt\Big]
\;=\;
\frac{\mathrm{Cov}^P\!\Big(M_T,\;\int_0^T\sigma_t^2\,dt\Big)}{E^P[M_T]},
\]
即 VRP 的符号由“随机贴现因子与已实现方差的协方差”决定。更直接的是：任何严格正的风险厌恶（即 $M_T$ 与方差冲击负向联动）都会使 $Q$ 测度给高波动态以更高权重，从而让 $E^Q[\sigma^2]>E^P[\sigma^2]$。也就是说，VRP 的符号与“投资者厌恶方差风险”这个前提一一对应：只要代表性投资者把高方差时期视作更差的状态（风险厌恶高、或下降端厌恶），$Q$ 就会系统地把方差抬价，卖出方差者便能持续收到正溢价。这解释了为何 VRP 与同期市场风险厌恶高度同向，也解释了为何央行压低风险厌恶时会同时压低 VRP。

\textbf{估计量之争：BTZ 与 Bekaert–Hoerova。} VRP 的实证测量并非唯一。BTZ（2009）用同期已实现方差 $RV_t$ 作为 $E^P[RV]$ 的替代，将其代换进 $VRP=VIX_t^2-RV_t$。这一做法在平稳期没有问题，但在剧烈波动期 $RV_t$ 会瞬时应激冲高，导致 $VRP$ 被强迫压成接近零甚至为负——从而制造出“危机中风险厌恶骤降”的假象，这与第 210 段“危机中恐惧反而回落”的度量问题同根。Bekaert 与 Hoerova（2014）改用状态空间/GARCH 类条件方差作为 $E^P[RV]$，用全样本日历信息而非同期实现值，得到的 VRP 序列更平滑、且在压力期不致反转。二者的分歧不在理论上，而在“现实测度期望”应由“同期已实现值”还是“模型条件期望”来代理——这也解释了为何不同论文对同一危机的 VRP 走势会给出相反读数。对读者的落点：报告 VRP 时必须同时给出测度口径，否则“VRP 在危机中显著转正/转负”这类陈述无从检验。

\textbf{上下行与偏度分解。} 把风险中性方差积分按执行价分成低于远期价（下行区间）与高于远期价（上行区间）两部分，得到下行方差风险溢价 DVP 与上行方差风险溢价 UVP。Londono 与 Xu（2019）证明，DVP 为正、逆周期、能较好预测四至六月收益，而 UVP 几乎只在极短期限有信号。机制上，DVP 的门槛在于“市场对左侧尾部（下跌）索取每单位方差更高的边际溢价”，UVP 则几乎不含系统性风险厌恶。因此从测量上，若只看总 VRP 会混入大量“不承载风险厌恶”的上行分量，削弱其作为情绪/风险厌恶指标的纯度——这是 VRP 应用中的一个隐蔽陷阱，也是文献用偏度而非方差作为“恐惧温度计”的原因之一。

对交易而言，卖 VIX 期货、卖方差互换或卖期权在约八成月份为正（业界常见口径 83\%–86\%，需人工复核），年均收取约 3–4 个波动点，形成正收益的 carry 型 alpha。但这并非免费午餐：正收益的前提是波动率不发生跳跃；一旦发生，负偏会把多年的盈利一次性吞没。因此对 VRP 的正收益，正确态度是意识到它同时携带“凸性脆弱”——平庸获利与尾部清零是同一枚硬币的两面。对宏观研究者，VRP 是衡量市场风险厌恶的、几乎实时且可高频计算的代理，被广泛用于货币政策与金融稳定研究。

局限方面：VRP 测度高度依赖条件方差如何估计（GARCH、已实现方差、HAR 等），不同模型在压力期甚至可给出方向相反的溢价序列（Chen 与 Cheng 2018 详细检视过这一点），因此“VRP 的位阶”不如“VRP 方向的一致性”稳健；此外，用同期已实现方差替代 P 测度期望的做法，在剧烈波动期会人为压低甚至逆转 VRP。

\begin{figure}[htbp]
\centering
\begin{tikzpicture}[scale=1.0,every node/.style={font=\small}]
% 总体柱表达 VIX^2（以平静期 VIX≈12 为例，VIX^2≈144 的方差尺度）
\draw[thick,thcolor,fill=blue!12] (0,0) rectangle (2.2,7.2)
   node[midway,align=left] {预期方差（现实测度 $E^P[RV]$）\\\footnotesize 平静期约对应 100–125\\ \scriptsize（VIX 约 10–11 的平方）};
% 上层为方差风险溢价（文献真实量级：长期为正，约 20–30 个方差点量级）
\draw[thick,htcolor,fill=teal!18] (0,7.2) rectangle (2.2,9.8)
   node[midway,align=center] {方差风险溢价（VRP）\\ \footnotesize 长期为正、约 20–30 方差点\\ \scriptsize VIX${}^2 - E^P[RV]$};
% 右侧总高标注
\draw[thick,->] (3.4,0)--(3.4,9.8)
   node[midway,right,align=left,text width=5.4cm] {总高对应 $VIX^2$\\（风险中性方差 $E^Q[RV]$，平静期约 144）};
% 右下真实量级说明框
\node[align=left,anchor=west,text width=6.8cm,font=\footnotesize] at (5.0,2.6) {
文献口径（需人工复核）：VRP 长期为正，\\业界常见为其中约 83–86\% 的月份为正，\\年均约 3–4 个波动点的水平。$\:\:$\\
即 $VIX > \sqrt{E^P[RV]}$ 通常成立。
};
\end{tikzpicture}
\caption{方差风险溢价（VRP）的分解与文献真实量级。将 $VIX^2$ 视为风险中性测度的期望方差，分离出“预期（现实）方差”与“方差风险溢价”两部分（Bekaert 与 Hoerova 2014）。量级依据文献通用口径：平静期 VIX 约 10–12、对应 $VIX^2$ 约 100–144；VRP 长期为正，约 83–86\% 的月份为正、年均约 3–4 个波动点（需人工复核）。图为示意结构，数值为文献量级而非逐日实测。}
\label{fig:vrp}
\end{figure}

% ==================================================================
% 第五章
% ==================================================================
\section{波动率期限结构与风险定价的斜率}

股票只有一条价格序列，但波动率市场却有一条完整的期限结构，包括 VIX 期货曲线以及从标普 500 不同期限期权中提取的 VIX1M、VIX3M、VIX6M 与 VIX1Y。期限结构的斜率、水平与形态是研究波动率预期的重要信号；更关键的是，它在多个期限上分别揭示了投资者如何为“较近的波动”与“较远的波动”定价。

近端与远端承载的信息不同。近月波动率（如 VIX9D）对短期新闻与恐慌高度敏感；远月（如 VIX3M、VIX1Y）更多反映市场对一个季度到一年状态的整体信念。斜率表现为升水（contango）或贴水（backwardation）：当近端高、远端低（贴水）时，市场在定价“当下紧张、未来趋平”，通常是恐慌形态；当近端低、远端高（升水）时，是平静中略带预期的形态。

这一领域近年最反直觉的发现来自 Dew-Becker、Giglio、Wang 与 Yaron（2017）。他们用方差互换曲线估计不同期限的方差风险溢价，发现只有“已实现波动率的冲击”承载显著的风险溢价；而对未来远期波动率的预期（forward volatility 的冲击），在超过约两个月之后几乎不再被定价——投资者愿为远期波动保险支付的溢价趋近于零，一个流传的说法是“超过两个月的对冲几乎是免费的”。这与期限结构向下倾斜的经验一致：Andries、Eisenbach、Kahn 与 Schmalz（2015/2025）用结构式与半参数两种方法，发现方差风险定价由 1–3 个月处最强，单调下降至 6–9 个月后几乎不显著，且压力环境下整条曲线更陡。Koijen 与 Van Nieuwerburgh 的讲义把这一现象直白表述为“波动长期保险近似免费”。

\textbf{仿射模型下的远期方差与期限结构。} 波动率的期限结构之所以可测、可建模，根基在于“远期方差”在一定模型族下可写成闭合形式的仿射函数。考虑 Heston 型随机波动率模型，期限 $T$ 的风险中性方差满足
\[
d\nu_t=\kappa_Q(\theta_Q-\nu_t)\,dt+\sigma_\nu\sqrt{\nu_t}\,dW^\nu_t,
\]
则风险中性测度下“从现在到 $T$ 的累计方差”的条件期望为
\[
E^Q_t\!\Big[\int_t^T \nu_s\,ds\Big]
\;=\;A(T-t)+B(T-t)\,\nu_t,
\]
其中 $A(\tau)=\theta_Q[\tau + (e^{-\kappa_Q\tau}-1)/\kappa_Q]$、$B(\tau)=(1-e^{-\kappa_Q\tau})/\kappa_Q$。沿同一逻辑可提取任意期限区间 $[t_j,t_{j+1}]$ 的\textbf{远期方差} $F_{t_j,t_{j+1}}=E^Q[\int_{t_j}^{t_{j+1}}\nu_s\,ds\mid\mathcal F_t]$，它在仿射设定下仍为 $\nu_t$ 的线性函数。把不同期限的 $F$ 连起来，就得到一条理论化的方差期限曲线——它的水平受 $\theta_Q$（长期均值）约束，斜率受 $\kappa_Q$（均值回归速度）驱动，曲率则来自 $\nu_t$ 相对 $\theta_Q$ 的偏离。这个结构解释了期限结构的核心事实：平静期 $\nu_t<\theta_Q$、曲线向上（升水），恐慌期 $\nu_t>\theta_Q$、曲线向下（贴水）。

\textbf{为什么远期波动“免费”：一个半群论证。} Dew-Becker 等（2017）的“远期波动近乎无溢价”并非反直觉的偶然，其根源在于波动率的均值回归与测度变换的相互作用。设想把“风险中性期望”拆成“物理期望 + 测度修正”：$E^Q_t[F]=E^P_t[F]+\Delta_t$。若波动率在 $P$ 下以高速 $\kappa_P$ 均值回归，则物理期望随期限延长的衰减极快；而测度修正项 $\Delta_t$ 取决于“方差风险容忍 λ 与剩余期限 $T-t$ 的乘积”。当 $T-t$ 拉长，波动的持续性早已消失，物理路径几乎回落到 $\theta_P$，此时即便 $\lambda$ 不为零，$E^Q-F$ 的差也会趋于平坦——即“远端几乎不再被定价”。换句话说，波动率的均值回归已经把长远期限的未来信息“洗白”了，只有短端仍保留被实时兑现的方差冲击。这正是半群性质导致的机制性结果，而非某种市场异常。

\textbf{对策略的含义（落地）。} 把“期限结构斜率”用作卖波动择时的信号，应同时看水平与斜率：当 $VIX$ 高位、曲线贴水（backwardation）、且 $VVIX/VIX$ 或隐含偏度异常抬升时，是系统性减仓卖空波动率的强信号；反之当 $VIX$ 低迷、曲线陡峭升水、且远期方差相对近端明显上升时，卖 VIX 期货的滚进收益较高，但须记住这不改凸性脆弱的本质——它只是提高了“收取保险费”的频率，并未消除“一次赔付多年”的尾部。实务中常以“近端—3 个月”或“曲线斜率”作为离散触发，但至少应在报告里标注口径（即用何种期限价差、何种去远期化方式），否则跨年比较会失真。

这一结论给“卖 VIX 期货拿到的是哪一种溢价”提供了精确答案：短端溢价来自已实现波动率冲击的跳跃风险，而非抽象的远期不确定性。对试图做远端方差套利的策略而言，这几乎是一盆冷水——远端几乎没有溢价可赚，近端溢价则是对真实跳跃的偿付。实务上，卖 VIX 期货的收益率主要由曲线升水时的滚进（roll-down）加上到期向现货收敛构成，这解释了为什么 VXX 这类跟踪长期波动率期货的指数天然存在每日负收益（常被称作磨损或衰减）。期限形态也可用于择时：当曲线进入贴水时，通常对应市场恐慌最盛、卖波动类策略最应减仓的时点；而当升水陡峭且现货低迷（如 2017 年）时，恰是系统化卖波动策略最赚钱也最拥挤的时段，这也是 2018 年 2 月事件的伏笔（见后文）。

局限在于稳健性与经济含义的测度依赖：条件方差与远期方差如何度量会影响斜率结论（需人工复核具体数值）；而“长期波动溢价趋近于零”这一结论建立在美国 1996–2014 年的样本上，跨期与跨市场的稳定性仍有争议。

\begin{figure}[htbp]
\centering
\begin{tikzpicture}
\begin{axis}[
    width=0.94\textwidth, height=7.2cm,
    xlabel={到期期限（月）}, ylabel={VIX 期货价},
    xmin=0, xmax=8.5, ymin=8, ymax=42,
    xtick={0,1,2,3,4,5,6}, ytick={10,15,20,25,30,35,40},
    grid=major, grid style={gray!25}, axis lines=left,
    legend style={at={(0.02,0.98)},anchor=north west,font=\footnotesize,draw=none},
]
% 现货点（真实区间锚点：平静期低点 9.14 位于 2017，示意近端）
\addplot[only marks,mark=square*,mark options={color=thcolor,scale=1.4}] coordinates {(0,16)};
\addlegendentry{现货 VIX（平静期约 12–17）}
% contango：升水向上（平静期，2017 观察到的形态）
\addplot[thick,thcolor,mark=*] coordinates {(0,16) (1,17.2) (2,18.1) (3,18.8) (4,19.3) (5,19.7) (6,20.0)};
\addlegendentry{升水（contango）}
% backwardation：贴水向下（恐慌期，近端高企）
\addplot[thick,ruled,mark=*,dashed] coordinates {(0,32) (1,29) (2,27) (3,25.6) (4,24.7) (5,24.1) (6,23.7)};
\addlegendentry{贴水（backwardation）}
% 真实要点框：置于右上部空白区（y>30 已超出曲线高度），远离虚线末端
\node[align=left,anchor=west,font=\footnotesize,fill=white,fill opacity=0.9,text opacity=1,draw=black!30] at (6.6,33) {真实要点\\· 2017-11-03 现货收 9.14（实测）\\· 2018、2024 恐慌时近端贴水\\· 曲线升水时卖出可“滚”得正收益};
% 简洁形态标注置于图内高处（y=39-40 空白区），远离曲线
\node[align=center,font=\footnotesize,fill=white,fill opacity=0.9,text opacity=1] at (4.2,39) {\scriptsize 贴水：近端高、远端低（恐慌期减仓）};
\node[align=center,font=\footnotesize,fill=white,fill opacity=0.9,text opacity=1] at (2.6,12) {\scriptsize 升水：远期高于近端（平静期滚进）};
\end{axis}
\end{tikzpicture}
\caption{VIX 期货期限结构的两种典型形态。曲线为结构示意：升水（实线）对应平静期，卖出方可获得滚进收益；贴水（虚线）对应恐慌期近端高企，是“卖波动”类策略的减仓信号。图中给出的是真实观测锚点：2017-11-03 现货收盘 9.14（实测），以及 2018、2024 年恐慌时近端贴水的形态特征。除标注为“实测”的数值外，曲线与坐标刻度属示意，不作精确引用。}
\label{fig:term}
\end{figure}

% ==================================================================
% 第六章
% ==================================================================
\section{衍生品定价与市场微观结构}

把 VIX 当作标的资产去定价其期货、期权与互换，会遇到一层二阶困难：波动率刻画的是价格变动的速度，而 VIX 期权是把这速度的变动速度也当作可交易对象。这一部分的难点在于，给定底层 SPX 期权市场的定价结构，VIX 及其派生品的价格如何形成，又在哪里与实际市场偏离。

技术上首先需要面对的，是 VIX 自身的分布形态。VIX 是标普 500 方差的 Q 测度期望再开平方，而 VIX 期货是未来某个 30 天隐含方差的期货；开方、到期与非线性共同决定了它呈高度单峰、负偏、尖峰且右尾极厚，任何假设其高斯的做法都会低估跳跃概率。VIX 期权的隐含波动率即 VVIX（波动率的波动率），刻画的是恐慌中的慌乱：当 VVIX 相对 VIX 的比值升高，意味着市场在给波动率本身的波动定价，通常对应系统性对冲需求激增。

定价框架方面，Heston 类随机波动模型及其加跳跃的扩展是常态工具；Andersen、Fusari 与 Todorov（2015，Econometrica）发展的从期权面板反演状态的非参数/参数混合方法，使“从 SPX 期权面板恢复 VIX 动力学”成为可能。方差互换定价与复制的理论基石是 Britten-Jones 与 Neuberger（2000）的静态复制条件，Carr 与 Wu（2009）则检验了其在真实市场中的“模型无关”边界。期限风险定价（承上一节）由 Dew-Becker 等（2017）与 Andries 等（2015/2025）给出：只有短端已实现方差冲击被定价，这为 VIX 衍生品的期限溢价划定了边界。微观结构侧的研究关注 VIX 期货上市后的价差、深度、未平仓量以及做市商再平衡行为；BIS 对 2018 年 2 月 5 日的拆解显示，VIX 衍生品市场的价格发现与风格再平衡相互咬合，日终再平衡可能在特定时点造成非线性的价格跳变（见后文）。

\textbf{一个实用定价框架：VIX 期货即“方差交换”，VIX 期权即“波动率交换的衍生”。} 记 $F_t=\mathbb E^Q_t[VIX_{T}]$ 为到期日 $T$ 的 VIX 期货价、$K$ 为执行价。若已能构建 VIX 的 $Q$ 测度密度，则欧式 VIX 期权价格可由折现期望 $e^{-r(T-t)}E^Q_t[(\max(VIX_T-K,0))]$ 给出，而这一期望的解析计算通常需要 VIX 自身的转移密度。在 Heston 框架下，$VIX_T$ 是 $V_T=\int_T^{T+30}\nu_s ds$ 的平方根的非线性函数，因而 VIX 期权的特征函数无法从 SPX 期权直接闭式得到——这正是“波动率的波动率”在建模上比普通股票期权难一层的根源。实践中常用两类替代：一是直接用 $V_T$ 的仿射矩做短期限近似（在 $T-t$ 较小时，$\sqrt{V_T}$ 的三阶展开足以抓住负偏与峰度）；二是借用 $VVIX$ 所隐含的 VIX 期权波动率，作局部随机波动率（local-stochastic vol）校准。无论哪类，都要记得底层是“车径上的车径”：任何对 SPX 期权尾部假设的错误，都会因二次取根被放大到 VIX 期权上。

\textbf{微观结构中的定价偏差机制。} 边际上，VIX 衍生品价格偏离理论值来自需求失衡而不是模型误差。当尾部对冲需求集中（如恐慌或事件前），做市商被迫在流动性稀薄时平仓，会出现“卖出方报价急速上移、深度收缩”的机械抬价；而当市场上出现大量单向 ETP（尤其反向波动 ETP），其日终再平衡形成可预测的机械购买力，常会在收盘窗口制造临时的价格气泡与随后回撤。这类偏差与“短线抢跑rebalance”同源（第十章展开）。对定价者而言，识别方式是检查“实际成交价 vs 理论模型价”的残差是否与 VVIX、做市商持仓、ETP 规模等拥挤度代理系统性地共动，而不是把它们当作偶然误差。

对应用的关键认识，一是持有 VIX 期货多头并不等于复制现货 VIX——二者分属不同期限、不同风险的合约，用它做尾部对冲时必须明确它保护的是未来 30 天隐含波动率的暴露，且随到期衰减滚动。二是对做市商而言，VVIX/VIX 之比以及 VIX 期权隐含的单日/单周尾部概率，是衡量波动率市场内部拥挤度的坐标，比值的异常抬升往往是系统性对冲需求的前兆。

VIX 衍生品定价模型的极限来自底层 SPX 期权面板的稀薄尾部——深度虚值期权往往缺报价——以及跨期拼接带来的噪声。任何模型在极端跳跃出现之前的表现都只是插值；“远、近”两层结构还会放大测量噪声，样本较小时结论并不稳健。

% ==================================================================
% 第七章
% ==================================================================
\section{交易策略的实证：卖波动率赚的是哪种钱}

这一章讨论的是最贴近实战、也最容易被夸大回报的部分：做空 VIX、卖出波动率或卖方差互换，究竟凭什么盈利，又在何种意义上只是“赚了报价的便宜”。现实中这类策略在样本内表现优异，却反复在实盘中被清算，理解这组矛盾是理解 VIX 资产属性的关键。

拆开来看，策略的收益来源可以归结为两个成分。一是 carry（代价率）：卖 VIX 期货的收益主要来自曲线升水时的滚进（roll-down）加上到期收敛，因曲线长期升水，卖出方能系统性取得持续为正的 carry。二是凸性脆弱（convex fragility）：这类策略的收益分布左偏且带极端负尾，绝大多数交易日小幅盈利（收取保险费），极小概率巨幅亏损（一次赔付多年的保费）。市场常以“在压路机前捡硬币”形容这种结构。因此，卖波动率的超额收益很难称之为真正的 alpha，它更接近为负偏收取的保费——其回报等价于做多偏度（skewness），而偏度在资产定价中长期被高额补偿，卖偏度者收取正收益、承担负偏。

各家族策略的实证可以分开看。其一，做空 VIX 期货或反向 VIX ETP，在 2010–2017 年这类区间内因波动率贴水加期限滚进而交出高 Sharpe（具体数值需人工复核），但 2018 年 2 月、2020 年 3 月与 2024 年 8 月的跳跃点会把多年收益一次清零。其二，卖出跨式/蝶式，同时赚时间衰减与卖偏度；Coval 与 Shumway（2001）、Bakshi 与 Kapadia（2003）的经典研究早已指出，平值跨式在无摩擦理想世界中赚取可观均值，但最大损失无界并承受巨大负偏。其三，Cboe 推出带防御的“卖波动”策略指数——VPN（卖出 VIX 期货但买入 VIX 看涨保护）与 VPD（动态去杠杆的 VIX 期货卖出）；Szado（2020）的测度显示这类设计通过在大跳跃前动态减杠杆或购买保护，显著收窄了最大回撤，为“既要 carry 又不愿被负偏清零”提供了模板（具体回撤数值需人工复核）。其四，波动率择时（Moreira 与 Muir 2017）把仓位与已实现波动率反比配置，样本内能显著提升 Sharpe，本质是以时间分散替代部分标的分散。

\textbf{偏度定价：卖波动率赚的正是“负偏的价格”。} 为了严谨说明“卖波动率≈卖偏度”，需要借助资产定价的普适公式：任意交易组合 $X$ 的无套利价格可用风险中性期望 $p(X)=E^Q[X]$ 表达；而把 $Q$ 相对 $P$ 的密度在 $X$ 的多项式上展开，可把风险中性期望写成现实期望加一系列“矩风险价格” $\times$ “$P$ 测度矩”的修正
\[
E^Q[X]\;\approx\;E^P[X]\;+\;(\text{mean price})\,E^P[X]+(\text{variance price})\,\mathrm{Var}^P[X]+(\text{skewness price})\,\mathrm{Skw}^P[X].
\]
对收益左偏的卖波动组合，$P$ 测度方差与其矩相对小，但\textbf{负偏（skewness）价格显著为正}——因为投资者在高偏的“灾难状态”里特别厌恶。于是，$E^Q[X]>E^P[X]$ 等价于“卖出波动率的现实收益 < 其应得风险补偿”，即卖家在平均意义上是亏的，只是把这部分“错价”当作 alpha 收取。这正是 Coval 与 Shumway（2001）、Bakshi 与 Kapadia（2003）在平值跨式上量化的对象。由此可得一个可检验推论：卖波动率的超额 Sharpe 应主要来自承担负偏敞口，而非其“扛波动”本身；若能剥离该敞口（如购买远端 VIX 看涨保护，即 VPN 思路），多余的正收益应显著减少。这为“带防御的卖波动”提供了理论正当性——它没有消除溢价，而是把负偏转嫁出去，从而让“carry 收益”与“尾部清零”解耦。

\textbf{波动率择时的机制：为什么“反比仓位”有效。} Moreira–Muir（2017）的策略是令组合权重 $w_t\propto 1/\sigma_t^2$。把它的收益分解，便可见其作用的要点：它并非“提高风险资产收益”，而是“降低持有期方差”。设单期收益为 $r_t=\mu+\varepsilon_t$、$\varepsilon_t$ 的方差 $\sigma_t^2$ 随机，按 $1/\sigma_t^2$ 配权后，组合波幅与标的波动率的乘积被压缩，而期望回报的衰减因子 $E[1/\sigma_t^2]$ 又比标的本身 $E[\sigma_t^2]$ 对状态不那么敏感——于是风险调整后收益（$E[r]/Std(r)$）上升。这也意味着该策略的增益对“波动率聚集”的强度、对执行噪声都敏感：在波动率平滑且可预测的环境中增益最大，在波动率呈跳跃时（$1/\sigma_t^2$ 对 $\sigma_t$ 的凸性被尾部放大）其平稳性减弱，样本内增益会明显下移。

\textbf{落地：把四种策略用法与风险控制串起来。} 对不同策略给出可执行但不构成建议的边界：
\begin{itemize}[leftmargin=2em,itemsep=2pt]
  \item 做空波动率类（卖 VIX 期货/卖方差互换/卖出跨式）：收益呈左偏厚尾，仓位应随 VIX 阈值与期限贴水动态缩量，且需预设“届时立即平仓”的硬止损映射到“一次性赔付多年”的尾部事件；不宜把历史高 Sharpe 当作常态。
  \item 带保护卖波动（VPN/VPD 型）：以买入远端 VIX 看涨或动态减杠杆对冲负偏，占用部分 carry 换取尾部保护；应把“保护成本”计入可承受的年度预算。
  \item 波动率择时（Moreira–Muir 型）：以滚动窗口的已实现方差或 VIX 反比仓，作为组合层面的“波动率目标”覆盖，通常放在组合层而非单个策略层实现。
  \item 用 VIX 做尾部对冲：明确其保护的是“未来 30 天隐含波动率”而非“指数点位本身”，需按“衰减曲线”滚动；近端贴水期成本最高，是减仓而非加仓时点。
\end{itemize}

实操中必须正视回测与现实之间的落差。许多回测使用日频或无费率模拟，忽略了跳跃日必须成交且支付点差的事实，加入更厚的尾部与更真实的执行成本后，卖波动策略的收益中枢会明显下移。拥挤本身也会放大负偏：当市场普遍做空波动率（2017、2024 年），一旦进入平仓，单一方向的拥挤会加速价格跳变。此外存在显著的生存偏差——能被长期观察到的卖波动产品，大多是清算事件中的幸存者，用幸存者绩效估计整体期望会系统性高估策略价值（样本筛选口径需人工复核）。

\begin{figure}[htbp]
\centering
\begin{tikzpicture}
\begin{axis}[
    width=0.9\textwidth, height=6.2cm,
    xlabel={策略收益（示意）}, ylabel={概率密度（示意）},
    xmin=-30, xmax=8, ymin=0, ymax=1.0,
    xtick={-30,-20,-10,0,8}, ytick distance=0.25,
    grid=major, grid style={gray!25}, axis lines=left,
]
% 左偏厚尾分布：主峰在略正收益，左尾延伸至大幅亏损
\addplot[thick,thcolor,smooth] coordinates {
    (-30,0.02) (-26,0.03) (-22,0.045) (-18,0.07) (-15,0.10) (-12,0.15)
    (-9,0.22) (-7,0.30) (-5,0.40) (-3,0.62) (-1,0.90) (0.5,1.0)
    (1.5,0.88) (2.5,0.72) (3.5,0.58) (4.5,0.44) (5.5,0.31) (6.5,0.18) (8,0.04)
};
% 均值线（略低于峰）
\draw[dashed,ruled] (axis cs:0,0)--(axis cs:0,1.0);
\node[fill=white,fill opacity=0.85,text opacity=1,anchor=west] at (axis cs:1.0,0.96) {\footnotesize 小正收益为主（多数月份赚钱）};
\node[align=center,fill=white,fill opacity=0.85,text opacity=1,anchor=east] at (axis cs:-16,0.16) {\footnotesize 极端负尾：\\ 一次跳跃吞掉多年收益};
\end{axis}
\end{tikzpicture}
\caption{卖波动率策略收益分布的“凸性脆弱”示意。收益集中在略为正的多数月份，而当 VIX 发生跳跃时左尾被拉成极端亏损；这就是“平庸获利、尾部清零”一体两面的来源。本图为示意性分布，不作精确概率引用。}
\label{fig:convex}
\end{figure}

% ==================================================================
% 第八章
% ==================================================================
\section{情绪与行为：恐慌指数究竟测了什么}

媒体把 VIX 称为恐慌指数，但心理学意义上的恐惧与资产定价意义上的波动率溢价并不等同。这一章回到行为金融的提问方式：VIX 走低或走高，究竟反映投资者对客观风险的理性评估，还是对坏结果的情绪性过度反应；它与“情绪”（sentiment）和“风险厌恶”（risk aversion）之间是怎样的关系。

可以把市场中的“恐惧”分解为三个理论上可分离的成分：其一，客观波动预期的抬高（物理测度下预期方差升高），属信息驱动；其二，风险厌恶的上升，即在同一预期方差下投资者要求更高溢价，属偏好驱动；其三，情绪或流动性驱动的过度反应，包括做市商再平衡与拥挤交易引发。这正是方差风险溢价 Q–P 分解在行为维度的翻版：VIX 平方相对已实现方差的超额部分，既可能是真实的“风险厌恶定价”，也可能是“情绪泡沫”。Bekaert 与 Hoerova（2014）的贡献之一，是主张用溢价成分捕捉风险厌恶、用条件方差成分捕捉对实体经济的担忧，为“VIX 里有多少情绪”提供了一个可操作的刻度。

\textbf{把“恐惧”做成可测的统计量。} 若将 VIX 平方 $E^Q[RV]$ 与用物理模型（如 GARCH/HAR）估计的 $E^P[RV]$ 都落实到日频，则它们的比值或差、以及期权的风险中性偏态（由虚值 put 相对 call 的定价差刻画），可以作为三个相对独立的“恐惧代理”：条件方差捕捉客观不确定性、VRP 捕捉风险厌恶、风险中性偏度捕捉左侧尾部情绪的强度。三者并不等比移动，因而可以搭建一个简单的判别矩阵：如果 VRP 与偏度同步抬升而条件方差平移不大，那么“恐惧”更多来自交易者风险厌恶与情绪；如果三项同步抬升，则更多反映客观风险恶化。这种拆分是 Q–P 主轴的直接应用，也为“VIX 高企到底是理性还是情绪”提供了可复用的操作而非口号。

\textbf{对策略的含义。} 若某段 VIX 的高企主要由情绪（VRP/偏度）驱动，那么它更可能在情绪回落后均值回归，卖波动类策略在“情绪高点”之后往往有较优的风险收益；反之若主要由客观条件方差驱动，则趋势更持久、快速回归的概率下降。因此，与其用单一 VIX 水平触发交易，不如把“VIX-条件方差差”与“偏度”作为二维触发，在区分情绪与基本面后再决定是顺应还是逆势。这同样是第四章 VRP 分解在情绪层的一种落地。

文献上，Whaley（2000）以“恐慌指标”命名并营销了这一概念，后续研究既沿用其作为指标，也批判其混淆波动率溢价与投资者恐惧。Baker–Wurgler 式情绪指数以及各类拥挤与尾部指标（如 VVIX/VIX 之比、VIX 期权隐含偏度）被用来定位 VIX 中的情绪成分；研究普遍发现，在恐慌期 VIX 抬升的同时，隐含偏度（put 相对 call 更贵）也同步走高——偏度是比波动率更纯粹的情绪信号。Bekaert、Hoerova 与 Lo Duca（2012）以 VAR 揭示宽松货币政策会压低 VRP 进而压低风险厌恶，提示卖波动率的部分收益其实是央行压低均衡风险厌恶的结果。行为一侧的证据则显示，VVIX 攀升与单一方向的拥挤交易（如长期做空波动的 ETP）会放大 VIX 的自我实现上冲；2018 年 2 月“白天几无动静、尾盘瞬间暴涨”的走势，就是情绪性或机制性过度反应的极端案例。

应用上，判断情绪不应只看 VIX 水平，而应把 VIX 与已实现波动率的差（溢价）放在一起看，更接近投资者愿意多付溢价的程度；对做市商与事件交易者，期权隐含偏度与尾部概率往往比 VIX 本身更早发出恐慌定价的信号；理解货币政策会临时压低风险厌恶，有助于解释为何卖波动在宽松期特别赚钱、在收紧期格外危险。

情绪并非一个可精确测度的单一变量，而是一簇代理的合集；把任何一个代理当作“真实情绪”，都隐含了强假设。风险厌恶与非理性情绪在时间序列上高度共线、难以干净分离，这里给出的只是刻度与方向，而非一刀切的判据。

% ==================================================================
% 第九章
% ==================================================================
\section{跨国与跨资产联动：作为风险传导节点的 VIX}

当美国 VIX 飙升时，新兴市场股票、原油、黄金乃至加密货币常常同步剧烈波动。这一现象引出两个问题：VIX 究竟是一个全球金融风险偏好的共同探测器，还是只是美国股市恐慌的出口；它对全球多资产组合管理者是否具有稳定的风险传导与领先意义。

金融市场的波动存在系统性共同成分。当单一市场（通常指美国）风险偏好收敛，全球投资者会同步调整头寸（portfolio rebalancing），波动随之传染到各资产，因此 VIX 更像是全球风险溢出网络中的一个中心节点，而非孤立的美股指标。美联储的研究（如 Londono 与 Wilson 2018）发现各国权益期权隐含波动率高度相关，若做市值加权会得到“全球隐含波动率”，其中美国因素与货币政策是主要驱动。使用 Diebold–Yilmaz 波动率溢出指数与 TVP-VAR 框架的近年研究（2018–2026 密集发表）普遍发现，VIX 是全球波动率的显著“净输出者”，传导具有方向性与不对称性：发达市场对其反应更强、更持久，新兴市场的暴露程度则各异（中国、俄罗斯、土耳其、墨西哥等长期同步性较强，部分市场只作阶段性暴露）；COVID 期间全球波动率系统的总关联明显抬升，疫情后回落（具体数值需人工复核）。跨资产层面，OVX（原油）、GVZ（黄金）以及加密货币波动指数与 VIX 呈正相关，但各资产自有基本面驱动，并非始终随美股同步。

\textbf{传导机制：为何 VIX 是“净输出者”。} VIX 之所以能作为全球波动的中心节点，其机制来自三道叠加的通道而非单一传染。第一，\textbf{共同宏观基石}：全球风险偏好的收敛（risk-off）本身是一个共同因子，任一市场的恐慌都会抬高全球同向定价，此时 VIX 既是共同因子的代理而非其唯一来源。第二，\textbf{套保—定价联动}：美股期权是全球机构对冲的基准标的，当美股 VIX 抬升，机构的组合对冲需求上升，会同时买入其他市场期权与波动率工具，从而把溢价“传导”给他处。第三，\textbf{流动性—去杠杆联动}：VIX 飙升期，做市商与风险平价基金为首要去杠杆对象，他们被迫在多个市场同步抛售，形成跨资产的机械式同步波动。正因为这三层，VIX 的溢出在危机期（三层同时激活）远强于平静期（仅第一层），这解释了为何总关联强度随 regime 大幅时变。也正因如此，用“VIX 涨→各市场跌”做风险模型时，不能只看线性相关，还要把它当成一个随状态切换强度、且可能被流动性冲击瞬间放大的结构。

\textbf{对组合的意义。} 对全球多资产组合，把 VIX 作为“统一风险闸门”是合理的：在 $VIX$ 与（VIX−条件方差）同时抬升、或隐含偏度恶化时，应统一收缩风险敞口，而不是按资产各自分散决策。但需注意，新兴市场对 VIX 的暴露并不均匀——长期同步性强的市场（中国、俄罗斯、土耳其等）会被更早按风险偏好定价，而部分市场仅阶段性暴露，因此"用 VIX 统一撤仓"的阈值应分市场差异化标定，避免对某些市场过早、却对另一些市场过晚撤出。

对多资产组合，VIX 飙升可作为撤敞口、备好流动性的系统性触发信号，常用于一体化的跨市场风控闸门。新兴市场与商品暴露在高 VIX 环境下往往与美股同步承压，理解这一点可以避免“以为分散、实则同涨同跌”。美联储的货币条件与全球波动率共同成分高度协同，提示货币政策是长期压低全球波动率的重要力量。

需要说明边界：全球共同成分并不等于“VIX 是因、其余是果”，溢出研究多建立在统计关联与预测误差分解之上，难以确立严格的因果方向；关联强度随市场状态高度时变，平静期低、危机期高，用历史均值估计未来传导强度会系统性偏低；各国 VIX 的构造口径也不尽一致，合并成“全球波动率”存在度量误差。

% ==================================================================
% 第十章
% ==================================================================
\section{极端事件与尾部风险：坠落中为何保险失效}

前面的各章多讨论波动率的共性规律，这一章转向毁灭时刻。VIX 的飙升往往极快，许多以 VIX 对冲的策略恰在最需要保护的下跌中失效甚至清零。理解 VIX 作为保险与作为赌具之间的窄缝，关键在于理解极端行情的形成机制。

极端行情下的 VIX 飙升很少单单由一条消息触发，更多是正反馈螺旋：价格下跌抬高已实现波动率，触发风控与去杠杆，进一步压低价格、再抬升波动率；与此同时，做市商或 ETP 发行人为维持日内目标敞口而做的机械式再平衡，会在尾部被迫接盘或抛售，把波动率的抬升放大成一个自洽的事实。因此 VIX 的极端形态几乎从来不是“平静后的自然上扬”，而是“短暂失速中的机械性坍塌与再膨胀”，这是它与普通资产回撤的重要差别。

几个关键事件可作档案。2008 年金融危机期间，VIX 于 11 月 20 日收报 80.86，盘中触及约 89.53（以 Cboe 口径为准），这场飙升与经济基本面同步，属客观波动与风险厌恶的叠加。2018 年 2 月 5 日，标普 500 当日仅跌约 4\%，VIX 却从 17.31 升到 37.32、涨幅 115.6\%，是自 1987 年崩盘以来最大的通胀率跳变；其机制被归因于拥挤的反向波动 ETP（XIV、SVXY 等）在日终被迫再平衡——一旦 VIX 上升、反向 ETP 必须买入 VIX 期货补空头损失，市场参与者提前抢跑，在尾盘把 VIX 期货价格推高成自我实现。BIS 测算当日 VIX 期货在 16:08 的一分钟内成交约 11.6 万手、约为全市场单日量的四分之一，价格被抬至远超预期信息的水平；结果是 XIV 触发加速清算（约损失 84\%–93\%）、SVXY 盘中熔断。2020 年 3 月 16 日，VIX 收报 82.69、刷新历史收盘高点并超过 2008 年，但在流动性救助下很快回落。2024 年 8 月 5 日因日元套息交易平仓，VIX 再次录得单日创纪录的飙升，说明单一拥挤交易的平仓即可在全球市场制造剧烈波动。

由此可提炼出几条教训。其一，做空波动的尾部不是“可能不会发生”，而是“何时发生”：只要 carry 足够诱人便会吸引资本，最终在某个拥挤日被自己的再平衡机制打穿，LJM、Catalyst、XIV 等被清算产品是同一剧本的不同演员。其二，用 VIX 对冲的成本与时机要求极高：平静期受滚进衰减拖累，极端期又因跳空而难以在最需要的位置成交；近端贴水恰表示市场已在恐慌中，此时买入的成本最高。其三，正反馈与挤兑是这一市场的结构性特征而非缺陷，识别拥挤度信号（曲线扁平化、反向 ETP 持仓激增、日终再平衡窗口的抢跑）比择时顶底更有意义。

\textbf{尾部建模：用极值理论而非正态插值。} 既然极端 VIX 事件的样本极少，经典正态或轻尾统计会在“百年一遇”级尾部严重低估概率。极值理论（EVT）提供了一条不预设分布形状的路径：对超阈值（POT）样本，渐近地，超出某一高阈值的超额分布可用广义帕累托分布（GPD）近似
\[
F_\xi(u)\;=\;1-\Big(1+\xi\,\frac{x-u}{\beta}\Big)^{-1/\xi},
\]
其中 $u$ 为阈值、$\beta$ 为尺度、$\xi$ 为形状参数。当 $\xi>0$ 时分布属 Frechet 型、尾部呈幂律厚尾，且 $\xi$ 越大尾部越厚。对 VIX 与已实现波动率的两类实测，形状参数估计通常落在正区间（典型约 $0.3$–$0.5$ 量级，需人工复核），这直接给出一个可检验推论：VIX 高尾出现的“可对冲频率”远高于“可被正态定价的频率”，因此“尾保险”的隐含成本必然系统性偏高。这一量化为“远月波动保险近乎免费，而短端跳跃保费昂贵”提供了极值层面的旁证：长端的均值回归已把尾部洗薄，短端却仍保留着最厚的、由真实跳跃撑起的尾部。

\textbf{为什么“最需要保险时买不到保险”。} 从正反馈机制可推出一个令人不快的结论：VIX 飙升时的“保险供给”会被抽干。因为尾部对冲的需求与 VIX 的抬升同向爆发，而做市商的库存与缓冲在极端行情下不足以同时满足所有买方；此刻卖方要么抬价到隐含概率大幅超过物理概率，要么干脆撤单。结果就是“保险价格在风险最高时达到最贵”，且成交深度骤降——这正是 2018 年 2 月 5 日晚间“价格在无成交量下被推高”的微观基础。对长期持有“买保险”策略的组合而言，这意味着两个实务含义：其一，保险应提前（在低波动期、曲线升水时）用较远月与分散到期滚动建仓，而非等到事件发生时临时购买；其二，评估对冲组合的价值不能只看其平均成本，还应计入“需要时可能无法按合理价成交”的流动性风险。也就是说，解决“坠落中保险失效”的办法并不是找到更好的尾部模型，而是改变建仓时点与执行方式本身——这是结构性问题，而非度量问题。

极端事件的样本极少，数十年中的真正断崖不过寥寥数次，任何从事件中归纳的模型都受小样本方差困扰，不能据此推断下一次将以何种形式、在哪个市场发生。此外，政策与流动性干预（如 2020 年美联储纾困）能迅速扭转 VIX 的形态，这说明极端事件的后验统计并不等同于无政策干预下的自然行为。

% ==================================================================
% 第十一章
% ==================================================================
\section{计量边界与机器学习时代}

近年的研究呈现出两个发展方向：机器学习在波动率预测上能走多远，以及哪些既有的 VIX 结论真正稳健、哪些只是特定样本与测量的产物。两者共同指向一个更根本的问题——波动率的可预测边界在哪里。

波动率的可预测性来自两个来源，也受一个限制约束。其一，波动率存在强烈的时变聚集（volatility clustering），今天的高波动往往预示明日的高波动，这是 GARCH、已实现方差、HAR 等模型都能捕捉的关系。其二，VIX 作为市场对未来的定价，天然携带前瞻信息。但极端跳跃（2008、2018、2020 年）本质上是低概率、高幅度的罕见事件，其精确时点几乎不可预测；模型能改进的是对风险权重或分布的估计，而非下一次冲击的日历。

计量上的挑战始终存在。条件的方差的测度方式会显著影响 VRP 等核心量的结论，任何对“VRP 显著”的断言都应做多种测度下的稳健性检验（Bekaert–Hoerova 2014 正是强调这一点）；2003 与 2014 两次方法变更使 VIX 序列存在非平稳段，跨断点的实证需谨慎分段；极端尾部与稀少的独立冲击使许多“显著”的因子在样本外难以复现，需警惕“本文显著、换种风格即消失”的多重检验风险。

机器学习方面，近年研究（如 Roszyk 与 Ślepaczuk 2024；Hortúa 与 Mora 2024）把 GARCH、LSTM、TCN、Transformer 及贝叶斯深度网络用于标普 500 与 VIX：将 VIX 纳入输入的混合 LSTM-GARCH 在样本内小幅优于纯 GARCH，概率深度学习（贝叶斯 TCN）在点预测与不确定性校准上表现出一定优势。这些进展更多是把已知的隐含信息与实际波动之间的拟合做细，而不是对未知跳跃的预言。另一个值得注意的进步是模型不再只给点预测，而是给出校准良好的概率区间，这对风险管理（而非择时）尤为有价值。需要提醒的是，机器学习论文常给人“波动率更可预测”的印象，但多数优势来自对历史轨迹的拟合，其真正的前瞻样外增益与可交易性仍需谨慎评估（具体指标需人工复核）。

综合而言，跨测度、跨样本相对稳健的结论包括：VIX 与收益的弱关联、VRP 长期为正并具一定收益预测力、短端方差风险定价强于长端、波动率聚集与均值回归。对模型与样本敏感的结论包括：具体的 VRP 数值、期限溢价下移的精确时点，以及各类卖波动策略的历史 Sharpe。下一次跳跃的时点与形态则几乎不可预测，任何声称能预测其时间的做法都应被高度怀疑。

% ==================================================================
% 结语
% ==================================================================
\section*{结语}
\addcontentsline{toc}{section}{结语}

把上述主题收拢起来看，VIX 研究给读者的核心教益是：金融市场为波动的不对称性定价，而 VIX 是把这种不对称性做成可交易商品的工具。它首先是测量工具，将预期波动与波动定价分离，为研究风险厌恶、情绪与恐慌提供实时坐标；其次是风险资产，拥有自身的期限结构、偏度与极端形态，定价难点在于那根厚重的右尾；再其次是风险源，作为全球波动溢出的中心节点，既检测风险也放大风险；最后是一台温度计，测出市场愿意为负偏支付的溢价，而人们往往在这一点上系统性犯错。

对 VIX 的预测力不应神化，对卖波动率的收益也不必妖魔化。更合适的做法是把它当作一台不对称性的定价仪器：用它理解风险、情绪与恐慌的分布；把它作为资产配置时，须记住买入的是保险或偏度，而这类保护的价格特性，往往在最不需要的时候最贵。

把各章收拢到开头提出的统一主轴上：几乎所有现象的根源都是 $Q$ 与 $P$ 的差距——它把“波动率”变成方差风险溢价（第四章）、把“期限”变成逐期限溢价的斜率（第五章）、把“偏度”变成负偏的价格（第六、七、八章）、把“尾部”变成厚尾与跳空（第十章）。我们反复看到四条根问题交错出现：哪些维度上 $Q$ 与 $P$ 分离（方差、偏度、厚尾都分离，而远期一维几乎不分离）；分离的幅度（可用偏度和 GPD 形状参数刻画）；为何持续存在（风险厌恶 + 保险需求 + 央行压低风险厌恶）；以及被何者放大（做市商再平衡、拥挤的 ETP、去杠杆）。这四问并不是并列的清单，而是一条因果链——测度分离保证了持续存在，偏度与厚尾放大了幅度，机械行为再把它瞬间引爆。理解这条链，就同时理解了“卖波动率长期为正”与“卖波动率偶尔清零”本是同一结构的两面。

给读者最终的判断：看待 VIX 的正确姿态，不是问“它值多少钱”或“下一次尖峰何时来”，而是持续问“此时 $Q$ 对哪一维度的定价，比 $P$ 所对应的可实现成本贵了多少、贵在哪个期限、由什么在放大”。把这三个问题的答案合起来，就构成一份可执行的 VIX 风险与策略坐标。参考文献按主题列于文末附录；正文中标“需人工复核”的数值建议在引用前以一手数据核验。

% ==================================================================
% 附录
% ==================================================================
\appendix
\section{关键文献（按主题选取）}

\begin{vtable}{>{\raggedright\arraybackslash}p{0.32\textwidth}>{\raggedright\arraybackslash}p{0.62\textwidth}}
Whaley (2000) & \textit{Stock Market Return and the Fear Gauge.} The Journal of Portfolio Management.\\
Bollerslev, Tauchen \& Zhou (2009) & \textit{Expected Stock Returns and Variance Risk Premia.} Review of Financial Studies, 22(11), 4463--4492.\\
Bekaert \& Hoerova (2014) & \textit{The VIX, the Variance Premium and Stock Market Volatility.} Journal of Financial Economics.\\
Bekaert, Hoerova \& Lo Duca (2012) & \textit{Risk, Uncertainty and Monetary Policy.} Journal of Monetary Economics.\\
Dew-Becker, Giglio, Wang \& Yaron (2017) & \textit{The Price of Variance Risk.} Journal of Financial Economics.\\
Cheng (2018) & \textit{The VIX Premium.} Review of Financial Studies, 32(1), 180--227.\\
Andries, Eisenbach, Kahn \& Schmalz (2015, 修订 2025) & \textit{The Term Structure of the Price of Variance Risk.} FRB NY Staff Reports 736.\\
Londono \& Xu (2019) & \textit{Variance Risk Premium Components and International Stock Return Predictability.} IFDP 1247.\\
Moreira \& Muir (2017) & \textit{Volatility-Managed Portfolios.} Journal of Finance.\\
Lochstoer \& Muir (2019) & \textit{Volatility Expectations and Returns.} NBER.\\
Coval \& Shumway (2001) & \textit{Expected Option Returns.} Journal of Finance.\\
Bakshi \& Kapadia (2003) & \textit{Delta-Hedged Gains and the Negative Market Volatility Risk Premium.} Review of Financial Studies.\\
Britten-Jones \& Neuberger (2000) & \textit{Option Prices, Implied Price Processes, and Stochastic Volatility.} Journal of Finance.\\
Carr \& Wu (2009) & \textit{Variance Risk Premiums.} Review of Financial Studies.\\
Andersen, Fusari \& Todorov (2015) & \textit{Parametric Inference and Dynamic State Recovery from Option Panels.} Econometrica.\\
Andersen, Bondarenko \& Gonzalez-Perez (2024) & \textit{VIX Maturity Interpolation.} Cboe Research.\\
Szado (2020) & \textit{Selling VIX Futures and Options for Portfolio Return Enhancement.} Providence College / Cboe.\\
Sushko \& Turner (2018) & \textit{The Equity Market Turbulence of 5 February: The Role of Exchange-Traded Volatility Products.} BIS Quarterly Review.\\
Diebold \& Yilmaz (2012) & \textit{Better to Give than to Receive: Predictive Directional Measurement of Volatility Spillovers.} International Journal of Forecasting.\\
Roszyk \& Ślepaczuk (2024) & \textit{The Hybrid Forecast of S\&P 500 Volatility Ensembled from VIX, GARCH and LSTM Models.} arXiv.\\
Hortúa \& Mora (2024) & \textit{Forecasting VIX Using Bayesian Deep Learning.} arXiv.\\
Heston (1993) & \textit{A Closed-Form Solution for Options with Stochastic Volatility.} Review of Financial Studies.\\
Bakshi, Kapadia \& Madan (2003) & \textit{Stock Return Characteristics, Skew Law, and the Differential Pricing of Individual Equity Options.} Review of Financial Studies.\\
Baker \& Wurgler (2006) & \textit{Investor Sentiment and the Cross-Section of Stock Returns.} Journal of Finance.\\
Pickands (1975) 及广义帕累托分布 & 极值理论（POT/GPD）阈值法的基础文献；亦见 Embrechts--Klüppelberg--Mikosch, \textit{Modelling Extremal Events}.\\
Zhang, Zhou \& Zhu (2009) & \textit{Excess Return of VIX Futures and Its Economic Value.} Journal of Financial Economics.（VIX 期货期限结构可用性的代表性研究）\\
\end{vtable}

\section{关键事件与数值对照年表}

下表的数值已经公开权威渠道核验，个别非官方口径以“约”字标注；正文中标记“需人工复核”的数值此处不再重复。

\begin{vtable}{l >{\raggedright\arraybackslash}p{0.30\textwidth} >{\raggedright\arraybackslash}p{0.48\textwidth}}
年份/日期 & 事件/数值 & 说明\\
1986--1989 & Brenner 与 Galai 提出“西格玛指数”思想 & VIX 的概念远祖，见《金融分析师期刊》。\\
1993 & Cboe 推出 VIX（初期基于标普 100 平值期权） & 最初由 Whaley 主持构造。\\
2003 & Cboe 与高盛合作更新 VIX 方法 & 改为基于标普 500、模型自由方差加权，回填至 1990；原标普 100 版更名 VXO。\\
2004-03-24 & VIX 期货在 CFE 上市 & 首个交易所交易的波动率期货。\\
2006-02 & VIX 期权在 C1 上市 & 现金结算、欧式。\\
2008-11-20 & VIX 收报 80.86，盘中约 89.53 & 金融危机期的历史高位（以 Cboe 口径为准）。\\
2014-10-06 & VIX 引入周度 SPX 期权 & 缩小到期期限插值误差约一个量级，形成统计断点。\\
2017-11-03 & VIX 收报 9.14（历史收盘低点） & 平静期极值。\\
2018-02-05 & VIX 单日涨 115.6\%（17.31→37.32） & 史上最大单日涨幅，触发 XIV 清盘。\\
2020-03-16 & VIX 收报 82.69（历史最高收盘） & COVID 冲击，超过 2008 年高点。\\
2024-08-05 & VIX 录得单日创纪录飙升 & 日元套息平仓驱动的全球波动巨震。\\
\end{vtable}

上述日期与数值仅作阅读上下文，不构成投资依据；官方最新历史数据以 Cboe 官网为准。

\section{术语速查}

\begin{vtable}{>{\raggedright\arraybackslash}p{0.25\textwidth}>{\raggedright\arraybackslash}p{0.69\textwidth}}
VIX & Cboe 标普 500 波动率指数，30 日模型自由隐含波动率。\\
VXO & 老版 VIX，基于标普 100 平值期权。\\
VVIX & “波动率的波动率”，基于 VIX 期权。\\
VIX9D / VIX3M 等 & 不同目标期限的 Cboe 波动率指数。\\
VRP & 方差风险溢价，Q 测度方差减 P 测度方差。\\
RV & 已实现（现实）波动率/方差。\\
Contango & 远期高于近端，期货升水（卖出方有利的结构）。\\
Backwardation & 近端高于远期，期货贴水（恐慌信号）。\\
Roll yield & 滚动换月的到期收敛收益，卖波动率 carry 的主要来源。\\
Carry & 持仓期内预期的静态与滚动收益之和。\\
Convex fragility & 凸性脆弱：收益左偏、尾部高爆，多数小赚、偶尔巨亏。\\
VPD / VPN & Cboe 动态去杠杆 / 带保护的两类“卖波动”策略指数。\\
BTZ & 指 Bollerslev--Tauchen--Zhou（2009）。\\
Q 测度 / P 测度 & 风险中性定价测度 / 现实概率测度。\\
\end{vtable}

\section{给研究者的自查清单}

无论是要复现文献结论，还是把 VIX 方法接到自己的项目，建议逐项自检。

\begin{enumerate}[leftmargin=2em,itemsep=2pt]
\item 口径是否一致：是否区分了现货 VIX、VIX 期货、VIX 期权与 VIX 曲线？三者对应不同的风险与费率，混用会使结论失真。
\item 结构断点是否处理：是否对 2003 与 2014 两个断点分段或加哑变量？否则跨断点回归可能被“变口径”误导。
\item VRP 的测度是否稳健：是否用 GARCH 与已实现/HAR 等两套口径交叉验证？单一测度下“VRP 显著”可能只是模型设定的产物。
\item 尾部与摩擦是否入账：是否计入跳跃日的执行成本、点差与拥挤折价？忽略它们会系统高估“卖波动”的样本内收益。
\item 是否被生存偏差误导：观察到的基本是清算事件中的幸存者，用幸存者平均估计整体期望会高估策略价值。
\item 显著性与经济含义是否分家：即便统计显著，若 $R^2$ 过小也不足以支撑实盘使用，小 $R^2$ 的预测并不等于可利用的 alpha。
\item 是否区分风险厌恶与情绪：高 VIX、高溢价既可能来自客观波动走高，也可能来自情绪挤兑。
\item 期限维度是否被忽视：远端波动保险相对便宜、短端才含跳跃溢价，避免在错误期限上追求 carry 或保护。
\item 全球关联是否按状态估计：波动关联在压力期骤升、平静期回落，用历史均值估计传导强度会有偏差。
\item 是否可复核：是否留下可追踪的数据源、模型版本与参数，使每个“需人工复核”之处可被验证。
\end{enumerate}

\vspace{1em}
\noindent\emph{补充说明：以上年份与出处以检索一致为准；个别具体数值（如逐月 VRP 均值、特定策略历史绩效）已在正文中以“需人工复核”标注，引用前请以一手数据源核验。}

\end{document}